\documentclass[a4paper,10pt]{report}\input{../0.Préambule/Préambule_cours_prof}\begin{document}

\pagestyle{empty}
\newpage
\section*{Pyramide et cône}
\addcontentsline{toc}{section}{Pyramide et cône}

\noindent \textit{L'usage de la calculatrice est \textbf{autorisée}.} \hfill \textit{Durée : 55 minutes}

\medskip
Un restaurateur, spécialiste dans la vente à emporter, cherche quel sera le meilleur emballage pour ses cornets de frites. Il a deux modèles à sa disposition :

\medskip
\begin{minipage}{3cm}
\begin{center}
\includegraphics[scale=0.2]{../40.Image/DS5 - 1.png}
\end{center}
\end{minipage}
\begin{minipage}{10cm}
Le modèle de gauche est une pyramide régulière
à base carrée de hauteur $24$cm et de largeur de
base $14$cm. (modèle P)

Le modèle de droite est un cône de révolution de
hauteur $24$cm et de rayon de base $8$cm.
(modèle C)
\end{minipage}
\begin{minipage}{3cm}
\begin{center}
\includegraphics[scale=0.2]{../40.Image/DS5 - 2.png}
\end{center}
\end{minipage}

\paragraph{Exercice 1}: \hfill \textbf{14 points \quad( Q1:2 / Q2:2 / Q3:3 / Q4:4 / Q5:3)}

\medskip
Se souvenant que la construction d'un cône de révolution est moins aisée, le restaurateur préfère étudier plus spécialement le cornet pyramidal.

\begin{enumerate}
\item Représenter en perspective la pyramide régulière.

Pour cela, on appelle : $ABCD$ sa base ; $S$ son sommet; $O$ le centre de sa base ; $I$ le milieu de l'arête $[BC]$

\vspace{3mm}
\item  Il commence à réaliser un patron de la pyramide en modèle réduis, \\mais ne sait plus comment le terminer.\\
Recopier le début de patron sur votre feuille puis compléter le patron \\de la pyramide sans oublier de mettre le nom de tous les points

\vspace{-2.3cm}
\begin{flushright}
\includegraphics[scale=0.4]{../40.Image/DS5 - 3.png}
\end{flushright}

\vspace{-5mm}
\item Quelle est la nature du triangle $OIS$ ? Justifier
\item  \textbf{Démontrer} que la longueur $SI$ mesure $25$cm. (on admettra que $OI = 7$cm)
\item  Pour connaître le coût de fabrication de son cornet, le restaurateur a besoin de connaître l aire de carton nécessaire à sa fabrication. Aide le à la déterminer.
\end{enumerate}

\paragraph{Exercice 2}: \hfill \textbf{6 points}

Pour se faire de la publicité, il affiche fièrement dans son restaurant sa participation au record du monde du plus grand cornet de frites datant du 22 septembre 2013.

\noindent \begin{minipage}{9cm}
\begin{doc}[Article de presse]
Le $22$ septembre $2013$, à Hazebrouck, $3 000$ éplucheurs
bénévoles pensent chiper aux Belges le record du plus
gros cornet de frites du monde.

Des milliers de personnes ont rempli
un cornet de $7,5$ mètres de haut

Vers $19$h, le compteur de la balance s'arrête et annonce
le poids du cornet à $6 225$ kilos. 

Le record de $1985$ est battu, largement, les Belges avaient réalisé un cornet de presque
cinq tonnes de frites. 

\end{doc}
\end{minipage}
\hspace{0.5cm}
\begin{minipage}{7.3cm}
\begin{doc}[Formules et données]

$\bullet$ $1$m$^3$ de pommes de terre pèse $622,5$kg

$\bullet$ $\text{Aire de la base} = \dfrac{3 \times \text{Volume} }{\text{Hauteur}} $

$\bullet$ $\text{Coté de la base} = \sqrt{\text{Aire de la base}}$

\end{doc}

\medskip
\begin{doc}[Le cornet]
Le cornet est assimilable à une pyramide régulière à base carrée.
\end{doc}
\end{minipage}

Déterminer la longueur du côté du carré de base du cornet.

\textit{Laisser apparent toute trace de recherche même non abouti}


\end{document}
